% id: 0165
% book: RPM 공통수학2 · 02 직선의 방정식
% section: 시험에 꼭 나오는 문제
% figure: none
% answer: ③
% answer_source: 계산(답 크롭 없음 · 해설 크롭 P(3, -1/3)까지 확인)
% variant_level: 0 (원본 전사)
% 이 파일은 build.mjs 가 items.json 에서 만든다. 고칠 때는 items.json 을 고친다.
\smash{\raisebox{1.5mm}{\dmkichul{RPM 공통수학2 0165번 · 28쪽 · 중요}}}
세 점~$\pt{A}(-1,\,1)$, $\pt{B}(5,\,-1)$, $\pt{C}(4,\,3)$을 꼭짓점으로 하는 삼각형~$\pt{ABC}$에서 선분~$\pt{AB}$ 위의 한 점~$\pt{P}$에 대하여 삼각형~$\pt{APC}$와 삼각형~$\pt{PBC}$의 넓이의 비가 $2:1$일 때, 두 점~$\pt{C}$, $\pt{P}$를 지나는 직선의 방정식은?
\begin{choicesii}{$3x-10y-18=0$}{$3x-10y+18=0$}{$10x-3y-31=0$}{$10x-3y+31=0$}{$10x+3y-49=0$}\end{choicesii}
